\noindent\textbf{(a)} En $x>0$, la rama del dato es $y=x/(1-\log x)$, $0<x<e$. Se prolonga por $x=0$ con $y(0)=y^{\prime}(0)=0$; a la izquierda se puede unir $y=0$ o una rama $y=x/(A-\log|x|)$, evitando su polo. La prolongación hacia $x<0$ no es única.

\noindent\textbf{(b)} $y=\pm\sqrt{2Cx+C^2}$, $C>0$, en $x>-C/2$, junto con $y=0$ en intervalos de $x>0$. No hay solución gráfica $C^1$ en un intervalo abierto que atraviese el origen.

\noindent\textbf{(c)} $y=\sqrt{x^2/2+7/(2x^2)}$, $x\in(0,\infty)$.

\noindent\textbf{(d)} $y=(3-x^2)/(2x)$, $x\in(0,\infty)$.
